% ============================================================================
%  Photon Shower Shape Cell-Level Reweighting: Methods, Implementation,
%  and Results
% ============================================================================
\documentclass[11pt,a4paper]{article}

% --- Packages ---------------------------------------------------------------
\usepackage[utf8]{inputenc}
\usepackage[T1]{fontenc}
\usepackage[margin=2.5cm]{geometry}
\usepackage{graphicx}
\usepackage{amsmath,amssymb}
\usepackage{booktabs}
\usepackage{longtable}
\usepackage{caption}
\usepackage{subcaption}
\usepackage{float}
\usepackage{xcolor}
\usepackage{hyperref}
\usepackage{listings}
\usepackage{array}
\usepackage{siunitx}

% --- Hyperref ---------------------------------------------------------------
\hypersetup{
  colorlinks=true,
  linkcolor=blue!50!black,
  citecolor=blue!50!black,
  urlcolor=blue!50!black
}

% --- Listings ---------------------------------------------------------------
\lstset{
  language=C++,
  basicstyle=\ttfamily\small,
  breaklines=true,
  frame=single,
  numbers=left,
  numberstyle=\tiny,
  keywordstyle=\color{blue!60!black},
  commentstyle=\color{green!40!black},
  stringstyle=\color{red!50!black},
  showstringspaces=false
}

% --- Graphics path ----------------------------------------------------------
\graphicspath{{/project/atlas/users/mfernand/QT/ShowerShapes/ShowerShapesAnalysis/output/data_mc_reweighting/combined/plots/}}

% --- Convenience macros -----------------------------------------------------
\newcommand{\Reta}{\ensuremath{R_{\eta}}}
\newcommand{\Rphi}{\ensuremath{R_{\phi}}}
\newcommand{\weta}{\ensuremath{w_{\eta 2}}}
\newcommand{\abseta}{\ensuremath{|\eta|}}
\newcommand{\chisq}{\ensuremath{\chi^2 / n_{\mathrm{df}}}}
\newcommand{\pt}{\ensuremath{p_{\mathrm{T}}}}

% --- Title ------------------------------------------------------------------
\title{\textbf{Photon Shower Shape Cell-Level Reweighting: Methods, Implementation, and Results}}
\author{Marta Fernandes da Silva}
\date{April 2026}

\begin{document}

\maketitle
\tableofcontents
\newpage

\begin{abstract}
This note documents the implementation and comparison of three cell-level
energy reweighting methods for correcting Monte Carlo simulation of photon
shower shapes (\Reta, \Rphi, \weta) in the ATLAS electromagnetic calorimeter
second layer. The methods operate on $7 \times 11$ cell clusters from
$Z \to \ell\ell\gamma$ events using Run 3 data (data24) and MC23e simulation.
Method 1 (flat shift) applies a constant per-cell correction. Method 2
(DeltaR-matched TProfile) follows Francisco Almeida's approach from
ATL-COM-PHYS-2021-640, using matched data--MC photon pairs to derive
fraction-dependent corrections. Method 3 (shift+stretch) matches both mean and
variance of each cell's energy fraction distribution. Results show that M1 and
M3 provide excellent agreement for most bins, while M2 suffers from systematic
energy-scale distortions caused by the matching strategy, degrading performance
especially for \weta. The ATLAS fudge factors remain the benchmark, with
\chisq values at the level of $\mathcal{O}(10^{-2})$ across the full detector.
\end{abstract}

\section{Introduction}

ATLAS applies shower-shape ``fudge factors'' to improve the agreement between
data and simulation for photon identification variables. These corrections are
very effective, but they act directly on the reconstructed shower-shape
observables. A cell-level method is conceptually more ambitious: if the energy
distribution in the underlying calorimeter cells is corrected, then all shower
shapes derived from those cells should improve simultaneously.

That is the motivation for this study. The target is the second layer of the
electromagnetic calorimeter, where the bulk of the photon energy is deposited
and where the three variables studied here are most sensitive:
\Reta, \Rphi, and \weta. The starting point is Francisco Almeida's Run 2 work
documented in ATL-COM-PHYS-2021-640. The present study ports that logic to Run
3 ntuples, validates it on $Z \to ee\gamma$, $Z \to \mu\mu\gamma$, and their
combination, and extends the comparison to a third method that performs the
actual shift-plus-stretch transformation at the cell level.

The work reported here should be read together with the broader Run 3 analysis
framework documented in ATL-COM-PHYS-2025-662. The emphasis of the present note
is technical and methodological: what was implemented, how it differs from the
reference Run 2 code, why Method 2 initially failed, why it still underperforms
after the software fixes, and what can realistically be improved next.

\section{Data and MC Samples}

The analysis uses Run 3 data24 photon ntuples and matching MC23e samples
produced with the NTupleMaker framework in Athena 25.0.40. Three channels are
processed:

\begin{itemize}
  \item $Z \to ee\gamma$ (``Zeeg''), built from the \texttt{egam3} data stream
  and the corresponding \texttt{mc\_eegamma.root} simulation.
  \item $Z \to \mu\mu\gamma$ (``Zmumug''), built from the \texttt{egam4} data
  stream and the corresponding \texttt{mc\_mumugamma.root} simulation.
  \item The combined sample, formed by merging both channels at histogram and
  correction level.
\end{itemize}

The ntuples are flat ROOT trees containing per-photon kinematics, the standard
ATLAS shower-shape branches (fudged and unfudged), and the full $7 \times 11$
Layer-2 cell energies and positions. The selected event yields after the full
analysis cuts are given in Table~\ref{tab:event-yields}.

\begin{table}[H]
\centering
\caption{Event selection summary after all cuts.}
\label{tab:event-yields}
\begin{tabular}{lrrr}
\toprule
Channel & Data events & MC events & Data/MC ratio \\
\midrule
$Z \to ee\gamma$ (Zeeg) & 56,344 & 810,895 & 0.069 \\
$Z \to \mu\mu\gamma$ (Zmumug) & 93,228 & 1,563,133 & 0.060 \\
Combined & 149,572 & 2,374,028 & 0.063 \\
\bottomrule
\end{tabular}
\end{table}

At the ntuple level, the raw trees contain approximately 8.9 million data events
for the $ee\gamma$ channel, 15.5 million data events for the $\mu\mu\gamma$
channel, and about 10 million plus 14.4 million MC events for the two signal
samples respectively. The cell-level analysis is therefore a strongly selected
FSR-photon measurement on top of much larger inclusive inputs.

\section{Event Selection}

The event selection is applied in two stages. The first stage is already
embedded in the NTupleMaker production; the second stage is applied inside the
reweighting macros in \texttt{scripts/data\_mc/config.h}.

\subsection{NTupleMaker-level preselection}

The ntuples are produced after the following baseline requirements:

\begin{itemize}
  \item Trigger, GRL, detector quality, and primary-vertex requirements.
  \item At least one photon and at least two leptons.
  \item Opposite-charge, trigger-matched leptons.
  \item Photon--lepton separation: $\Delta R(\gamma, e) > 0.2$ for electrons
  and $\Delta R(\gamma, \mu) > 0.01$ for muons.
  \item Broad radiative-$Z$ windows: $40 < m_{\ell\ell} < 120$~GeV and
  $45 < m_{\ell\ell\gamma} < 125$~GeV.
\end{itemize}

\subsection{Offline selection applied in the reweighting pipeline}

The reweighting analysis then tightens the selection in order to isolate a
clean FSR-photon sample and to enforce a well-defined cell topology:

\begin{itemize}
  \item FSR mass windows: $40 < m_{\ell\ell} < 83$~GeV and
  $80 < m_{\ell\ell\gamma} < 100$~GeV.
  \item Photon transverse momentum: $\pt > 7$~GeV.
  \item Photon pseudorapidity acceptance: $\abseta < 2.37$, excluding the crack
  region $1.37 < \abseta < 1.52$.
  \item Tight photon identification: \texttt{tightOffInc}.
  \item Unconverted photons only.
  \item Photon--lepton overlap removal: $\Delta R(\gamma, \ell) > 0.4$.
  \item Lepton--lepton separation: $\Delta R(\ell, \ell) > 0.2$.
  \item Exactly 77 Layer-2 cells in the cluster.
  \item The central cell ($k = 38$ in 0-based indexing) must be the hottest
  cell in the cluster.
  \item MC only: truth matching is required.
  \item MC event weights are the product of the generator weight and the pileup
  reweighting weight, although the M1/M2/M3 population moments are accumulated
  unweighted, as discussed below.
\end{itemize}

The detector is divided into 14 standard pseudorapidity bins,
\begin{equation}
  \{0, 0.2, 0.4, 0.6, 0.8, 1.0, 1.2, 1.3, 1.37, 1.52, 1.6, 1.80, 2.0, 2.2, 2.4\},
\end{equation}
with bin 8 corresponding to the crack region and therefore always empty.

\begin{table}[H]
\centering
\caption{Summary of the offline analysis requirements applied in the cell-level reweighting pipeline.}
\label{tab:selection-summary}
\begin{tabular}{p{5.4cm}p{8.5cm}}
\toprule
Requirement & Applied cut \\
\midrule
Mass window & $40 < m_{\ell\ell} < 83$~GeV, $80 < m_{\ell\ell\gamma} < 100$~GeV \\
Photon kinematics & $\pt > 7$~GeV, $\abseta < 2.37$, crack veto \\
Photon quality & tight ID, unconverted only \\
Angular separation & $\Delta R(\gamma,\ell) > 0.4$, $\Delta R(\ell,\ell) > 0.2$ \\
Cluster topology & 77 cells, central cell hottest \\
MC-specific & truth matched \\
\bottomrule
\end{tabular}
\end{table}

\section{Shower Shape Variables and Validation Strategy}

Three shower-shape variables are studied throughout this note:

\begin{align}
  \Reta &= \frac{E(3 \times 7)}{E(7 \times 7)}, \\
  \Rphi &= \frac{E(3 \times 3)}{E(3 \times 7)}, \\
  \weta &= \sqrt{\frac{\sum_i E_i \eta_i^2}{\sum_i E_i} -
                   \left(\frac{\sum_i E_i \eta_i}{\sum_i E_i}\right)^2}.
\end{align}

Here $E(7 \times 7)$ denotes the energy in the central $7 \times 7$ subset of
the Layer-2 cluster, $E(3 \times 7)$ the energy in the central three rows and
seven columns, and $E(3 \times 3)$ the energy in the central $3 \times 3$
core. The width variable \weta is computed in a $3 \times 5$ window using the
cell positions, consistently with the ATLAS implementation.

Two parallel histogram sets are built in the validation macro:

\begin{itemize}
  \item \textbf{Set A: branch-based.} The values are read directly from the
  ntuple branches \texttt{photon.reta}, \texttt{photon.rphi},
  \texttt{photon.weta2}, and their \texttt{unfudged\_*} versions.
  \item \textbf{Set B: cell-computed.} The values are recomputed from the
  stored $7 \times 11$ cell energies using the same positional windows as the
  cell-reweighting code.
\end{itemize}

The distinction matters. The branch-stored and cell-computed quantities are not
identical event by event: a small systematic offset of approximately
$0.006$--$0.008$ is observed, caused by differences between the branch-level
calculation and the explicit position-based reconstruction in the validation
macros. This offset is not a bug in the reweighting itself; it is a definition
difference. For that reason, Set A and Set B are compared separately throughout
this note.

\section{ATLAS Fudge Factors as Benchmark}

The standard ATLAS fudge factors provide the reference performance level. They
are applied at reconstruction level as direct shifts of the MC shower-shape
variables in bins of detector region and kinematics. In the present samples,
the fudged distributions show excellent data--MC agreement:

\begin{itemize}
  \item \Reta: typically $0.00$--$0.07$ in \chisq.
  \item \Rphi: typically $0.00$--$0.04$ in \chisq.
  \item \weta: typically $0.00$--$0.02$ in the barrel and only slightly larger
  in the endcap.
\end{itemize}

This is the target that a cell-level method should aim to match. Any new method
must therefore be judged against two baselines: the unfudged MC and the already
very strong fudge-factor performance.

\section{Method 1 --- Flat Shift}

Method 1 corrects only the mean of each cell-energy-fraction distribution. Let

\begin{equation}
  f_k = \frac{E_k}{E_{\mathrm{tot}}}
\end{equation}

be the energy fraction carried by cell $k$. The per-cell correction is

\begin{equation}
  \Delta_k = \langle f_{\mathrm{data}} \rangle_k - \langle f_{\mathrm{MC}} \rangle_k,
\end{equation}

and the corrected MC energy is

\begin{equation}
  E'_k = E_k^{\mathrm{MC}} + \Delta_k \cdot E_{\mathrm{tot}}^{\mathrm{MC}}.
\end{equation}

Because the energy fractions satisfy $\sum_k f_k = 1$ in both data and MC, the
sum of the shifts obeys

\begin{equation}
  \sum_k \Delta_k = 0,
\end{equation}

so Method 1 preserves the total cluster energy by construction and does not
need any renormalization step.

The core accumulation in \texttt{derive\_corrections.C} is

\begin{lstlisting}
// Accumulate weighted fraction for each cell
for (int k = 0; k < kClusterSize; ++k) {
    double frac = cellE->at(k) / Etot;
    stats.sumF[k][etaBin] += w * frac;
    stats.sumF2[k][etaBin] += w * frac * frac;
}
stats.sumW[etaBin] += w;
\end{lstlisting}

and the application in \texttt{validate\_data\_mc.C} is simply

\begin{lstlisting}
std::vector<double> m1_cells(kClusterSize);
for (int k = 0; k < kClusterSize; ++k) {
    m1_cells[k] = cellE->at(k) + m1_delta[k][etaBin] * Etot;
}
\end{lstlisting}

In the current implementation, these population moments are accumulated
unweighted. This is deliberate. Using the very large raw MC event weights in
the cell-population averages would corrupt the intended estimate of the
population mean and is also not how Francisco's original correction code was
used.

Method 1 is robust, simple, and numerically stable. Its limitation is equally
clear: it moves only the mean, not the width or higher moments of the cell
fraction distributions. This is particularly relevant for \weta, which is
sensitive to the spread of the energy profile in $\eta$.

\begin{table}[H]
\centering
\caption{Method 1 checksums for the combined channel. The sum of the per-cell corrections is consistent with zero in every pseudorapidity bin, confirming exact energy preservation.}
\label{tab:m1-checksum}
\begin{tabular}{crrc}
\toprule
Eta bin & Data events & MC events & $\sum_k \Delta_k$ \\
\midrule
0 & 17807 & 283121 & $6.6 \times 10^{-16}$ \\
1 & 21137 & 344005 & $-6.9 \times 10^{-15}$ \\
2 & 20156 & 329922 & $1.4 \times 10^{-14}$ \\
3 & 21820 & 344890 & $-8.9 \times 10^{-15}$ \\
4 & 16379 & 271357 & $3.4 \times 10^{-15}$ \\
5 & 13916 & 225763 & $4.0 \times 10^{-15}$ \\
6 & 6480 & 102928 & $-2.1 \times 10^{-15}$ \\
7 & 1530 & 24820 & $-2.7 \times 10^{-15}$ \\
8 & 0 & 0 & 0 \\
9 & 3437 & 46749 & $1.7 \times 10^{-15}$ \\
10 & 7889 & 117913 & $6.1 \times 10^{-15}$ \\
11 & 7673 & 113773 & $5.1 \times 10^{-15}$ \\
12 & 7224 & 106818 & $-3.2 \times 10^{-15}$ \\
13 & 4124 & 61969 & $-4.3 \times 10^{-16}$ \\
\bottomrule
\end{tabular}
\end{table}

\section{Method 2 --- DeltaR-Matched TProfile}

Method 2 is the most involved and also the most instructive part of this study.
It is the direct Run 3 adaptation of the profiled correction strategy described
in ATL-COM-PHYS-2021-640.

\subsection{Origin in Francisco Almeida's Run 2 workflow}

Francisco's original implementation consists of two main stages.
\texttt{NTUP.cxx} first builds a matched data--MC tree by pairing photons in
$(\eta,\phi)$ space. Then \texttt{tree.C} fills per-cell TProfiles that store
the average correction as a function of the MC cell energy fraction. The logic
is that a cell with a very small fraction may need a different correction than a
cell with a large fraction, so a flat shift is too restrictive.

The matching in \texttt{NTUP.cxx} is a nearest-neighbour search within
$\Delta R < 0.1$:

\begin{lstlisting}
for(int i = 0; i<MC_vect.size(); i++) {
    double drMin = 100;
    int k = -99;
    for(int j = 0; j<Data_vect.size(); j++) {
        double dr = MC_vect[i].first.DeltaR(Data_vect[j].first);
        if (dr > 0.1) continue;
        if (dr < drMin){
            drMin = dr;
            k = j;
        }
    }
    if(k == -99) continue;
    std::pair<int , int > test(MC_vect[i].second,Data_vect[k].second);
    i_j.push_back(test);
}
\end{lstlisting}

The corresponding TProfile binning in \texttt{tree.h} is

\begin{lstlisting}
const float EBins [12 + 1] =
{-1., -0.5, -0.4, -0.3, -0.2, -0.1, 0., 0.1, 0.2, 0.3, 0.4, 0.5, 1.};
\end{lstlisting}

This full $[-1,1]$ range is important because noise-subtracted cells can carry
negative energies and therefore negative fractions.

\subsection{Run 3 implementation}

Francisco's Run 2 ntuples and branches are not directly compatible with the Run
3 trees used in the present analysis. The method therefore had to be ported.
Instead of consuming a pre-matched tree, the Run 3 implementation reads the
data and MC trees independently, stores the selected photons in memory, and
performs the matching on the fly.

The matching algorithm implemented in \texttt{derive\_corrections.C} is:

\begin{enumerate}
  \item Read all selected data photons and store $(\eta,\phi)$, total energy,
  eta bin, and full 77-cell energies.
  \item Read all selected MC photons and store the same information.
  \item For each MC photon, find the nearest data photon in $(\eta,\phi)$.
  \item Require $\Delta R < 0.1$.
  \item If several MC photons point to the same data photon, keep only the
  closest pair.
\end{enumerate}

The naive implementation is $\mathcal{O}(N_{\mathrm{data}} N_{\mathrm{MC}})$,
which is intractable for the combined sample:

\begin{equation}
  149{,}572 \times 2{,}374{,}028 \approx 3.55 \times 10^{11}
\end{equation}

pair comparisons. The code was therefore optimized with an $\eta$-grid spatial
index. Data photons are binned in grid cells of width 0.1 in $\eta$, and each
MC photon only searches the corresponding cell and its two neighbours.

\begin{lstlisting}
const double kEtaGridLo = -3.0, kEtaGridHi = 3.0;
const double kGridCell = kMatchDR;
const int kNGridBins = (int)((kEtaGridHi - kEtaGridLo) / kGridCell) + 1;
std::vector<std::vector<int>> etaGrid(kNGridBins);

for (int id = 0; id < (int)dataPhotons.size(); ++id) {
    int bin = (int)((dataPhotons[id].eta - kEtaGridLo) / kGridCell);
    if (bin >= 0 && bin < kNGridBins)
        etaGrid[bin].push_back(id);
}

for (int imc = 0; imc < (int)mcPhotons.size(); ++imc) {
    int etaBin0 = (int)((mcPhotons[imc].eta - kEtaGridLo) / kGridCell);
    double bestDR = 999;
    int bestData = -1;
    for (int db = -1; db <= 1; ++db) {
        int gbin = etaBin0 + db;
        if (gbin < 0 || gbin >= kNGridBins) continue;
        for (int id : etaGrid[gbin]) {
            // compute dR and keep best match
        }
    }
}
\end{lstlisting}

This reduces the runtime from effectively unusable to practical for the full
combined sample.

\subsection{Definition of the profiled correction}

For each matched pair and for each cell $k$, the correction stored in the
TProfile is

\begin{equation}
  \alpha_k(f_{\mathrm{MC}}) = f_{k,\mathrm{data}} - f_{k,\mathrm{MC}},
\end{equation}

filled as a function of the MC cell fraction $f_{k,\mathrm{MC}}$. The TProfile
therefore stores the average additive correction required for a cell that enters
the shower with a given fraction of the cluster energy.

The application step is

\begin{equation}
  E'_k = E_k^{\mathrm{MC}} + \alpha_k(f_k^{\mathrm{MC}}) \cdot E_{\mathrm{tot}}^{\mathrm{MC}},
\end{equation}

followed by an explicit renormalization,

\begin{equation}
  E''_k = E'_k \cdot \frac{E_{\mathrm{tot}}^{\mathrm{MC}}}{\sum_j E'_j},
\end{equation}

because the sum of the profiled corrections is not constrained to vanish. If a
TProfile bin has fewer than 10 entries, the code falls back to the M1 flat
shift for that cell and eta bin.

\begin{lstlisting}
if (m2_prof[k][etaBin]) {
    int bin = m2_prof[k][etaBin]->FindBin(mc_frac);
    if (m2_prof[k][etaBin]->GetBinEntries(bin) >= 10)
        alpha = m2_prof[k][etaBin]->Interpolate(mc_frac);
    else
        alpha = m1_delta[k][etaBin];
}
m2_cells[k] = cellE->at(k) + alpha * Etot;
\end{lstlisting}

\subsection{Differences from Francisco's original code}

\begin{table}[H]
\centering
\caption{Detailed comparison between Francisco Almeida's original Run 2 implementation and the Run 3 implementation used in this study.}
\label{tab:francisco-comparison}
\begin{tabular}{p{3.2cm}p{5.4cm}p{5.4cm}}
\toprule
Aspect & Francisco (\texttt{NTUP.cxx} + \texttt{tree.C}) & Run 3 implementation \\
\midrule
Matching input & Pre-matched \texttt{Tree\_matched.root} & Matching performed on the fly from separate data and MC files \\
Branch names & Run 2 style branches (\texttt{MC\_ph\_*}) & Run 3 ntuple branches (\texttt{photon.pt}, \texttt{photon.eta}, etc.) \\
TProfile binning & $[-1,1]$ with 12 non-uniform bins & Same after the bug fix \\
Profile content & $\alpha = f_{\mathrm{data}} - f_{\mathrm{MC}}$ & Same \\
Application formula & \texttt{NewMC[k] += alpha} & $E'_k = E_k + \alpha \cdot E_{\mathrm{tot}}$ \\
Renormalization & Explicit rescaling after application & Same \\
Selection & Run 2 cuts, healthy clusters & Run 3 FSR selection, 77 cells, central hottest, truth matching \\
Performance scaling & Brute-force matching & Eta-grid spatial index \\
\bottomrule
\end{tabular}
\end{table}

\subsection{The dimensional bug in Francisco's application step}

One of the most important findings of the code comparison is that Francisco's
\texttt{tree.C} applies the correction as

\begin{lstlisting}
float alpha = CellCorrection[k][n][1]->Interpolate(
    MC_ph_clusterCellsLr2E7x11->at(k) / sumE_MC);
NewMC->at(k) = NewMC->at(k) + alpha;
\end{lstlisting}

This is dimensionally inconsistent: $\alpha$ is a fraction, while
\texttt{NewMC->at(k)} is an energy. The correct implementation is to convert the
fraction shift into an energy shift by multiplying by the total cluster energy.
The Run 3 port fixes this. This matters conceptually even if the original Run 2
workflow happened to produce tolerable numerical performance in the restricted
phase space used there.

\subsection{The historical TProfile binning bug}

The first Run 3 implementation of Method 2 used a TProfile range of $[0,0.5]$
with 9 bins. That choice was wrong. A large fraction of the cell distribution,
especially for peripheral and noise-subtracted cells, lies below zero. Those
entries were therefore pushed into underflow and effectively lost from the
profile. This was the main reason the earliest Method 2 results were
catastrophic, with \chisq values in the range 4--25 across many bins.

Replacing the bins by Francisco's original non-uniform
$\{-1,-0.5,-0.4,\ldots,0.5,1\}$ edges fixed that software error and improved
the Method 2 behaviour substantially, but it did not make the method
competitive with M1 or M3. That remaining gap is physical and methodological,
not just a coding bug.

\subsection{Why Method 2 still underperforms after the software fixes}

The smoking gun is the renormalization factor. If Method 2 produced a balanced
set of cell corrections, the sum of the corrected energies would already be very
close to the original cluster energy, and the renormalization factor would sit
near unity. That is not what is observed. For the combined sample, the average
Method 2 scale factor sits around 0.91, and drops to 0.87 in bin 9. In other
words, the raw Method 2 correction overshoots the cluster energy by roughly
7--13\% before the final rescaling.

\begin{table}[H]
\centering
\caption{Average renormalization scale factors for Methods 2 and 3 in the combined channel. Method 2 requires a large energy rescaling, while Method 3 stays extremely close to unity.}
\label{tab:renorm-scales}
\begin{tabular}{ccc}
\toprule
Eta bin & M2 $\langle E_{\mathrm{tot}} / \sum E'_k \rangle$ & M3 $\langle E_{\mathrm{tot}} / \sum E'_k \rangle$ \\
\midrule
0 & 0.930895 & 1.000032 \\
1 & 0.928092 & 1.000060 \\
2 & 0.922837 & 1.000027 \\
3 & 0.914850 & 1.000064 \\
4 & 0.920981 & 1.000056 \\
5 & 0.911042 & 1.006321 \\
6 & 0.918304 & 1.000065 \\
7 & 0.914107 & 1.000037 \\
9 & 0.871503 & 1.013186 \\
10 & 0.908250 & 1.000048 \\
11 & 0.930886 & 1.000024 \\
12 & 0.941118 & 1.000034 \\
13 & 0.916202 & 1.000033 \\
\bottomrule
\end{tabular}
\end{table}

This is the core problem. The $\Delta R$ matching pairs photons that are close
in detector coordinates, but not necessarily close in actual shower properties.
In a separate data/MC matching problem, there is no true event-by-event
correspondence. The TProfile therefore learns not only the average data--MC
difference, but also noise from mismatched shower topologies. Summed over 77
cells, that noise develops a positive bias and injects too much energy into the
cluster. The mandatory renormalization then scales every cell back down, which
distorts the relative shape and damages exactly the variable that is most width
sensitive: \weta.

This means that the remaining underperformance of Method 2 is not evidence of a
remaining software bug in the Run 3 code. It is evidence that the matching idea
itself is too weak when implemented as a nearest-neighbour search in
$(\eta,\phi)$ across independent data and MC samples.

\subsection{Matched-pair statistics}

The matching efficiency is high in every non-empty eta bin. For the combined
sample, 149,019 data--MC pairs survive the deduplication step, corresponding to
99.6\% of the selected data photons.

\begin{table}[H]
\centering
\caption{Method 2 matched-pair statistics for the combined channel. The final column shows the mean absolute correction for the central cell ($k=38$).}
\label{tab:m2-matched-combined}
\begin{tabular}{crc}
\toprule
Eta bin & Matched pairs & $\langle |\alpha| \rangle$ for cell 38 \\
\midrule
0 & 17,737 & 0.062358 \\
1 & 21,070 & 0.061408 \\
2 & 20,088 & 0.085019 \\
3 & 21,737 & 0.076187 \\
4 & 16,323 & 0.071177 \\
5 & 13,889 & 0.064589 \\
6 & 6,443 & 0.087800 \\
7 & 1,508 & 0.043685 \\
9 & 3,424 & 0.067682 \\
10 & 7,865 & 0.048136 \\
11 & 7,639 & 0.047263 \\
12 & 7,176 & 0.054153 \\
13 & 4,120 & 0.060074 \\
\bottomrule
\end{tabular}
\end{table}

\section{Method 3 --- Shift+Stretch}

Method 3 implements the actual shift-plus-stretch transformation at the cell
level. For each cell $k$, the corrected fraction is defined as

\begin{equation}
  f'_k = \frac{\sigma_{k,\mathrm{data}}}{\sigma_{k,\mathrm{MC}}}
         \left(f_k^{\mathrm{MC}} - \mu_{k,\mathrm{MC}}\right)
         + \mu_{k,\mathrm{data}}.
\end{equation}

The corresponding corrected energy is then

\begin{equation}
  E'_k = f'_k \cdot E_{\mathrm{tot}}^{\mathrm{MC}}.
\end{equation}

This is the ``real'' shift+stretch method in the sense that both the mean and
the variance of the cell-fraction distribution are explicitly matched. That is
precisely why Method 3 improves variables like \weta better than Method 1,
which only moves the mean.

The application step in \texttt{validate\_data\_mc.C} is

\begin{lstlisting}
double mc_frac = cellE->at(k) / Etot;
double f_new = m3_sigRatio[k][etaBin]
             * (mc_frac - m3_meanMC[k][etaBin])
             + m3_meanData[k][etaBin];
m3_cells[k] = f_new * Etot;
\end{lstlisting}

Method 3 still requires a final renormalization because independent cell-wise
stretching does not guarantee that the corrected fractions sum to unity.
However, unlike Method 2, the correction is naturally balanced and the
renormalization factor stays extremely close to 1.0, typically at the level of
$10^{-5}$--$10^{-3}$.

The practical complication is the sigma ratio
$\sigma_{\mathrm{data}} / \sigma_{\mathrm{MC}}$. In low-statistics bins,
especially near the barrel--endcap transition, this ratio can become extremely
large. The implementation therefore caps it to the interval $[0.5, 2.0]$.
That cap is necessary to avoid unphysical behaviour, but it also truncates the
correction exactly in the bins where the width mismatch is most problematic.
This is why Method 3 remains excellent overall but still shows visible failures
in eta bins 5 and 9.

\section{Results: Combined Channel}

\subsection{Set A: branch-based reference and fudge-factor benchmark}

The Set A comparison uses the shower-shape branches stored directly in the
ntuples. It therefore provides the cleanest reference for the ATLAS fudge
factors.

\begin{table}[H]
\centering
\caption{Combined-channel Set A \chisq values for branch-stored shower shapes. ``unf'' denotes the unfudged MC branches and ``fud'' the standard ATLAS fudged branches.}
\label{tab:setA-combined}
\begin{tabular}{crrr}
\toprule
Eta bin & \Reta (unf, fud) & \Rphi (unf, fud) & \weta (unf, fud) \\
\midrule
0  & 0.01, 0.00 & 0.04, 0.00 & 0.06, 0.00 \\
1  & 0.01, 0.00 & 0.03, 0.00 & 0.08, 0.00 \\
2  & 0.01, 0.00 & 0.05, 0.00 & 0.10, 0.00 \\
3  & 0.02, 0.00 & 0.02, 0.00 & 0.10, 0.00 \\
4  & 0.03, 0.00 & 0.07, 0.01 & 0.03, 0.00 \\
5  & 0.02, 0.00 & 0.09, 0.01 & 0.03, 0.00 \\
6  & 0.03, 0.01 & 0.07, 0.01 & 0.07, 0.01 \\
7  & 0.19, 0.07 & 0.06, 0.04 & 0.05, 0.01 \\
9  & 0.07, 0.03 & 0.02, 0.00 & 0.67, 0.01 \\
10 & 0.13, 0.01 & 0.02, 0.01 & 0.82, 0.02 \\
11 & 0.19, 0.01 & 0.11, 0.01 & 1.22, 0.01 \\
12 & 0.37, 0.01 & 0.03, 0.01 & 2.04, 0.01 \\
13 & 0.53, 0.03 & 0.02, 0.01 & 2.36, 0.00 \\
\bottomrule
\end{tabular}
\end{table}

The pattern is clear. The unfudged MC already models \Reta and \Rphi fairly
well, but \weta degrades strongly toward the endcap. The fudge factors recover
excellent agreement everywhere. This establishes the scale of the problem and
the benchmark that the cell-level methods have to match.

\subsection{Set B: cell-computed results for M1, M2, and M3}

The Set B comparison uses shower shapes recomputed from the $7 \times 11$ cell
energies after applying the cell-level corrections.

\begin{table}[H]
\centering
\caption{Combined-channel Set B \chisq values for cell-computed shower shapes.}
\label{tab:setB-combined}
\begin{tabular}{crrr}
\toprule
Eta bin & \Reta (MC, M1, M2, M3) & \Rphi (MC, M1, M2, M3) & \weta (MC, M1, M2, M3) \\
\midrule
0  & 0.01, 0.01, 0.14, 0.01 & 0.04, 0.00, 0.29, 0.00 & 0.06, 0.01, 1.56, 0.02 \\
1  & 0.01, 0.01, 0.08, 0.00 & 0.03, 0.00, 0.27, 0.00 & 0.08, 0.01, 2.90, 0.01 \\
2  & 0.01, 0.01, 0.08, 0.01 & 0.05, 0.00, 0.31, 0.00 & 0.09, 0.02, 3.61, 0.02 \\
3  & 0.02, 0.00, 0.10, 0.00 & 0.02, 0.01, 0.58, 0.01 & 0.10, 0.00, 1.97, 0.01 \\
4  & 0.03, 0.01, 0.18, 0.02 & 0.07, 0.01, 0.72, 0.01 & 0.03, 0.01, 3.25, 0.02 \\
5  & 0.02, 0.01, 0.17, 0.09 & 0.09, 0.01, 0.63, 0.24 & 0.03, 0.01, 5.89, 3.33 \\
6  & 0.03, 0.02, 0.24, 0.03 & 0.07, 0.01, 0.77, 0.03 & 0.07, 0.01, 2.20, 0.01 \\
7  & 0.16, 0.12, 0.56, 0.15 & 0.06, 0.06, 0.75, 0.06 & 0.04, 0.02, 1.13, 0.01 \\
9  & 0.07, 0.37, 0.07, 0.33 & 0.02, 0.24, 0.10, 0.42 & 0.67, 0.01, 1.91, 2.93 \\
10 & 0.13, 0.01, 0.39, 0.02 & 0.02, 0.02, 0.91, 0.02 & 0.82, 0.01, 4.68, 0.01 \\
11 & 0.19, 0.07, 10.83, 0.08 & 0.09, 0.04, 7.65, 0.08 & 1.22, 0.01, 24.10, 0.01 \\
12 & 0.37, 0.03, 1.44, 0.03 & 0.03, 0.02, 4.65, 0.01 & 2.06, 0.02, 6.19, 0.00 \\
13 & 0.53, 0.05, 2.16, 0.04 & 0.02, 0.02, 24.40, 0.01 & 2.37, 0.01, 1.56, 0.00 \\
\bottomrule
\end{tabular}
\end{table}

Several conclusions follow immediately.

First, the cell-computed baseline is consistent with the branch-based baseline:
the absolute values are close, confirming that the Set B reconstruction from
cells is valid.

Second, Method 1 works remarkably well. It improves the baseline almost
everywhere and brings \Reta and \Rphi to very small \chisq values over most of
the detector. It also strongly improves \weta in the endcap, although that
variable remains sensitive to width effects that M1 cannot fully address.

Third, Method 2 is unstable. In a few bins it helps or remains comparable to the
baseline, but in several bins it is catastrophically worse. The clearest cases
are eta bin 11 for all three variables and eta bin 13 for \Rphi.

Fourth, Method 3 is the best cell-level method overall. In the barrel it is at
the level of M1 or better, and for \weta it usually outperforms M1 because it
acts on the variance as well as the mean. Its failures are localized, not
global, and are directly traceable to sigma-ratio capping in bins 5 and 9.

\subsection{Average performance summary}

Table~\ref{tab:average-chi2} summarizes the average performance over all
non-empty eta bins in the combined channel.

\begin{table}[H]
\centering
\caption{Average \chisq across all non-empty eta bins in the combined channel. The averages are computed directly from Tables~\ref{tab:setA-combined} and \ref{tab:setB-combined}.}
\label{tab:average-chi2}
\begin{tabular}{lccc}
\toprule
Method & $\Reta$ & $\Rphi$ & $\weta$ \\
\midrule
MC baseline & 0.12 & 0.05 & 0.65 \\
Fudge factors & 0.01 & 0.01 & 0.01 \\
M1 (flat shift) & 0.06 & 0.03 & 0.01 \\
M2 (TProfile) & 1.26 & 3.23 & 4.69 \\
M3 (shift+stretch) & 0.06 & 0.07 & 0.49 \\
\bottomrule
\end{tabular}
\end{table}

The hierarchy is unambiguous. The fudge factors remain the best overall
correction. Among the cell-level methods, M1 is the most robust and M3 is the
most complete physically. M2 is not competitive in its present form.

\section{Per-Channel Results}

The Zeeg and Zmumug channels show the same qualitative behaviour as the
combined sample, which is important because it shows that the main conclusions
are not an artefact of the channel combination.

\begin{table}[H]
\centering
\caption{Zeeg Set B \chisq values for cell-computed shower shapes.}
\label{tab:setB-zeeg}
\begin{tabular}{crrr}
\toprule
Eta bin & \Reta (MC, M1, M2, M3) & \Rphi (MC, M1, M2, M3) & \weta (MC, M1, M2, M3) \\
\midrule
0  & 0.02, 0.02, 0.14, 0.02 & 0.05, 0.01, 0.26, 0.01 & 0.06, 0.01, 1.87, 0.02 \\
1  & 0.01, 0.01, 0.09, 0.03 & 0.04, 0.01, 0.27, 0.01 & 0.08, 0.01, 3.83, 0.01 \\
2  & 0.01, 0.02, 0.07, 0.01 & 0.04, 0.00, 0.26, 0.00 & 0.09, 0.03, 5.05, 0.03 \\
3  & 0.03, 0.01, 0.07, 0.01 & 0.03, 0.01, 0.50, 0.01 & 0.10, 0.01, 2.47, 0.01 \\
4  & 0.03, 0.01, 0.13, 0.02 & 0.08, 0.02, 0.64, 0.02 & 0.03, 0.03, 2.66, 0.03 \\
5  & 0.03, 0.03, 0.13, 0.07 & 0.11, 0.05, 0.56, 0.28 & 0.03, 0.03, 12.99, 4.54 \\
6  & 0.07, 0.09, 0.24, 0.07 & 0.07, 0.02, 0.89, 0.06 & 0.09, 0.02, 3.72, 0.01 \\
7  & 0.19, 0.16, 0.53, 0.19 & 0.09, 0.10, 1.00, 0.09 & 0.08, 0.04, 1.21, 0.09 \\
9  & 0.07, 1.22, 0.28, 0.30 & 0.04, 1.30, 0.08, 0.48 & 0.55, 0.07, 2.58, 5.50 \\
10 & 0.13, 0.03, 0.30, 0.02 & 0.07, 0.05, 0.87, 0.08 & 0.69, 0.02, 5.57, 0.04 \\
11 & 0.23, 0.09, 7.62, 0.09 & 0.17, 0.14, 4.02, 0.04 & 1.12, 0.03, 17.60, 0.02 \\
12 & 0.36, 0.02, 1.24, 0.03 & 0.02, 0.01, 3.60, 0.02 & 1.71, 0.01, 5.59, 0.01 \\
13 & 0.60, 0.10, 1.56, 0.10 & 0.04, 0.02, 3.88, 0.02 & 2.23, 0.03, 1.35, 0.02 \\
\bottomrule
\end{tabular}
\end{table}

\begin{table}[H]
\centering
\caption{Zmumug Set B \chisq values for cell-computed shower shapes.}
\label{tab:setB-zmumug}
\begin{tabular}{crrr}
\toprule
Eta bin & \Reta (MC, M1, M2, M3) & \Rphi (MC, M1, M2, M3) & \weta (MC, M1, M2, M3) \\
\midrule
0  & 0.01, 0.01, 0.15, 0.01 & 0.04, 0.00, 0.32, 0.00 & 0.07, 0.01, 1.87, 0.02 \\
1  & 0.01, 0.01, 0.09, 0.01 & 0.03, 0.00, 0.29, 0.00 & 0.07, 0.01, 3.39, 0.01 \\
2  & 0.02, 0.01, 0.10, 0.01 & 0.05, 0.00, 0.35, 0.00 & 0.10, 0.01, 3.30, 0.01 \\
3  & 0.02, 0.00, 0.11, 0.01 & 0.03, 0.01, 0.61, 0.01 & 0.10, 0.00, 2.50, 0.01 \\
4  & 0.03, 0.01, 0.19, 0.02 & 0.07, 0.01, 0.74, 0.02 & 0.03, 0.01, 3.48, 0.02 \\
5  & 0.02, 0.01, 0.19, 0.01 & 0.09, 0.01, 0.76, 0.01 & 0.04, 0.01, 6.70, 0.01 \\
6  & 0.02, 0.01, 0.27, 0.02 & 0.09, 0.02, 0.81, 0.03 & 0.07, 0.01, 2.36, 0.01 \\
7  & 0.22, 0.18, 1.12, 0.15 & 0.09, 0.08, 1.06, 0.07 & 0.06, 0.05, 1.94, 0.02 \\
9  & 0.11, 0.10, 0.22, 0.21 & 0.02, 0.02, 0.59, 0.21 & 0.77, 0.02, 3.21, 0.04 \\
10 & 0.14, 0.01, 0.36, 0.03 & 0.02, 0.02, 0.83, 0.03 & 0.94, 0.01, 18.02, 0.01 \\
11 & 0.19, 0.07, 9.92, 0.14 & 0.07, 0.01, 2.59, 0.10 & 1.28, 0.02, 11.89, 0.01 \\
12 & 0.40, 0.05, 1.74, 0.05 & 0.04, 0.02, 5.49, 0.02 & 2.51, 0.03, 5.33, 0.01 \\
13 & 0.58, 0.06, 2.65, 0.07 & 0.03, 0.02, 27.23, 0.02 & 2.50, 0.04, 1.67, 0.00 \\
\bottomrule
\end{tabular}
\end{table}

The more statistical Zmumug sample gives slightly smoother behaviour in some
regions, but the pathology of Method 2 is common to both channels. The same is
true for the difficult eta bins in Method 3. This reinforces the interpretation
that the observed issues are intrinsic to the correction strategies, not a
channel-specific fluctuation.

\subsection{Problematic bins}

Three patterns deserve explicit emphasis.

\begin{itemize}
  \item \textbf{Method 1:} the largest residuals appear in bins 9 and 11--13
  for \weta, where matching the mean alone is not sufficient to correct the
  width of the lateral profile.
  \item \textbf{Method 2:} eta bins 11 and 13 are catastrophic. The worst case
  is \Rphi in bin 13 with \chisq $= 24.40$, followed by \weta in bin 11 with
  \chisq $= 24.10$.
  \item \textbf{Method 3:} bins 5 and 9 are the problematic regions, exactly
  where the sigma-ratio capping becomes active because of unstable variance
  estimates near the barrel--endcap transition.
\end{itemize}

\section{Plots}

This section collects the full set of combined-channel plots produced by the
pipeline. All figures use the combined sample because it provides the most
stable visual summary; the channel-split tables above demonstrate that the same
conclusions hold separately in Zeeg and Zmumug.

\subsection{Branch-based shower-shape comparison}

\begin{figure}[H]
\centering
\includegraphics[width=0.92\textwidth]{branch_reta.pdf}
\caption{Distribution of \Reta for data (black), unfudged MC (blue), and fudged MC (red) in all non-empty eta bins. The fudge factors bring the MC into excellent agreement with data across the full detector.}
\label{fig:branch-reta}
\end{figure}

\begin{figure}[H]
\centering
\includegraphics[width=0.92\textwidth]{branch_rphi.pdf}
\caption{Distribution of \Rphi for data (black), unfudged MC (blue), and fudged MC (red). The improvement is strongest in the forward bins, where the unfudged simulation shows clear residual shape differences.}
\label{fig:branch-rphi}
\end{figure}

\begin{figure}[H]
\centering
\includegraphics[width=0.92\textwidth]{branch_weta2.pdf}
\caption{Distribution of \weta for data (black), unfudged MC (blue), and fudged MC (red). This is the hardest observable, and the endcap mismatch of the unfudged MC is immediately visible. The fudge factors largely remove it.}
\label{fig:branch-weta2}
\end{figure}

\subsection{Cell-computed reweighting comparison}

\begin{figure}[H]
\centering
\includegraphics[width=0.92\textwidth]{rew_reta.pdf}
\caption{Cell-computed \Reta distributions for data (black), baseline MC (blue), Method 1 (green), Method 2 (red), and Method 3 (purple). Methods 1 and 3 track the data well; Method 2 remains visibly unstable in the forward bins.}
\label{fig:rew-reta}
\end{figure}

\begin{figure}[H]
\centering
\includegraphics[width=0.92\textwidth]{rew_rphi.pdf}
\caption{Cell-computed \Rphi distributions after the three cell-level methods. The catastrophic failure of Method 2 in bin 13 is visible as a gross shape distortion, whereas Methods 1 and 3 remain close to the data.}
\label{fig:rew-rphi}
\end{figure}

\begin{figure}[H]
\centering
\includegraphics[width=0.92\textwidth]{rew_weta2.pdf}
\caption{Cell-computed \weta distributions after the three cell-level methods. Method 3 is the most successful because it modifies the variance as well as the mean. Method 2 shows broad residual distortions caused by its large renormalization.}
\label{fig:rew-weta2}
\end{figure}

\subsection{Cell profile maps before and after correction}

\begin{figure}[H]
\centering
\includegraphics[width=0.92\textwidth]{cell_profiles_before.pdf}
\caption{Mean $7 \times 11$ cell-energy-fraction maps before any correction. The left panel shows data and the right panel shows the raw MC. The dominant differences are localized in the shower core and in the near-peripheral cells that drive \weta.}
\label{fig:cell-before}
\end{figure}

\begin{figure}[H]
\centering
\includegraphics[width=0.92\textwidth]{cell_profiles_after_m1.pdf}
\caption{Mean $7 \times 11$ cell maps after Method 1. The average data profile is reproduced almost exactly, as expected from the construction of the flat-shift correction.}
\label{fig:cell-after-m1}
\end{figure}

\begin{figure}[H]
\centering
\includegraphics[width=0.92\textwidth]{cell_profiles_after_m2.pdf}
\caption{Mean $7 \times 11$ cell maps after Method 2. The average profile is improved in some bins, but systematic residual patterns remain because the renormalization step distorts the intended cell-by-cell shifts.}
\label{fig:cell-after-m2}
\end{figure}

\begin{figure}[H]
\centering
\includegraphics[width=0.92\textwidth]{cell_profiles_after_m3.pdf}
\caption{Mean $7 \times 11$ cell maps after Method 3. The agreement is generally very strong, with the main residuals concentrated in the bins where sigma-ratio capping is active.}
\label{fig:cell-after-m3}
\end{figure}

\begin{figure}[H]
\centering
\includegraphics[width=0.92\textwidth]{cell_profiles_delta.pdf}
\caption{Data--MC difference maps before correction. These residuals define the pattern that the three correction methods are trying to remove. The most visible discrepancies are not limited to the central cell and therefore justify a genuine cell-level correction strategy.}
\label{fig:cell-delta-before}
\end{figure}

\subsection{Residual maps after correction}

\begin{figure}[H]
\centering
\includegraphics[width=0.92\textwidth]{cell_delta_m1.pdf}
\caption{Residual map after Method 1. The mean cell fractions are matched by construction, so the remaining visible structures mainly reflect the fact that the distributions are not Gaussian and that a mean-only correction cannot reproduce all higher moments.}
\label{fig:delta-m1}
\end{figure}

\begin{figure}[H]
\centering
\includegraphics[width=0.92\textwidth]{cell_delta_m2.pdf}
\caption{Residual map after Method 2. Coherent residual patterns remain, demonstrating that the profile-based correction plus renormalization does not preserve the intended shape transformation.}
\label{fig:delta-m2}
\end{figure}

\begin{figure}[H]
\centering
\includegraphics[width=0.92\textwidth]{cell_delta_m3.pdf}
\caption{Residual map after Method 3. The remaining discrepancies are localized and broadly consistent with the bins affected by sigma-ratio capping.}
\label{fig:delta-m3}
\end{figure}

\subsection{Stored-versus-computed validation}

\begin{figure}[H]
\centering
\includegraphics[width=0.92\textwidth]{stored_vs_computed_reta.pdf}
\caption{Comparison of stored and cell-computed \Reta values. The small offset confirms that the two implementations are not exactly identical, which is why branch-based and cell-based validations are kept separate.}
\label{fig:stored-computed-reta}
\end{figure}

\begin{figure}[H]
\centering
\includegraphics[width=0.92\textwidth]{stored_vs_computed_rphi.pdf}
\caption{Comparison of stored and cell-computed \Rphi values. The agreement is tight but not exact, consistent with small differences in the window definitions and positional treatment.}
\label{fig:stored-computed-rphi}
\end{figure}

\begin{figure}[H]
\centering
\includegraphics[width=0.92\textwidth]{stored_vs_computed_weta2.pdf}
\caption{Comparison of stored and cell-computed \weta values. This is the most sensitive observable to implementation details, and the offset observed here motivates the use of the dedicated Set B histograms for the cell-level comparison.}
\label{fig:stored-computed-weta2}
\end{figure}

\section{Discussion and Possible Improvements}

\subsection{Why Method 2 does not work well enough}

Method 2 was expected to perform better than Method 1 because it is more
flexible: instead of applying a single correction per cell and eta bin, it uses
an energy-fraction-dependent correction derived from matched data--MC pairs.
After debugging the binning and fixing the application formula, that expectation
is not borne out by the data.

The reason is that the matching procedure is not physically selective enough.
Two photons close in $(\eta,\phi)$ are not necessarily equivalent in cluster
topology, total energy, upstream material history, or leakage pattern. The
profiled correction therefore learns a mixture of true data--MC difference and
pairing noise. Summed over 77 cells, that noise generates a net positive energy
bias, and the compulsory renormalization then washes out the intended
cell-by-cell transformation.

Several improvements are plausible:

\begin{itemize}
  \item Match in a more physical space, for example $(\pt, \abseta)$ or a
  multidimensional feature set rather than only $(\eta,\phi)$.
  \item Constrain the profiled corrections such that
  $\sum_k \alpha_k(f_k^{\mathrm{MC}}) = 0$ on average, removing the large energy
  bias at the source.
  \item Apply a damped correction, for example using a factor smaller than one,
  to control the overshoot and study whether the current profile is simply too
  aggressive.
  \item Replace pairwise matching altogether by conditional response modelling
  in bins of the relevant photon properties.
\end{itemize}

\subsection{Why Method 3 still fails in bins 5 and 9}

The Method 3 failures are localized and understandable. Bins 5 and 9 sit near
regions where the calorimeter geometry and shower development change more
rapidly. In those bins, the observed variance ratios become unstable, with raw
values reaching well beyond any physically sensible range. The cap to $[0.5,2]$
prevents nonsense corrections, but once a large fraction of cells are clipped,
the transformation is no longer the exact data-matching shift-plus-stretch that
Method 3 intends to implement.

Possible improvements include combining neighbouring bins for the variance
estimate, introducing a smoother regularization of the sigma ratio rather than a
hard cap, or adding an explicit dependence on \pt so that broad detector bins do
not mix photons with substantially different shower topologies.

\subsection{Why Method 1 underperforms for weta2 in some bins}

Method 1 is mean-correcting by construction. That is ideal for observables whose
disagreement is dominated by a shift in the central tendency. It is not enough
when the disagreement is driven by the width of the distribution. Since \weta is
itself a width observable, Method 1 cannot be expected to solve all of its
residuals in bins where the data and MC variances differ substantially.

\subsection{General directions for future work}

The next steps suggested by this study are:

\begin{itemize}
  \item introduce explicit \pt-dependent corrections;
  \item explore multivariate cell corrections that preserve inter-cell
  correlations rather than treating each cell independently;
  \item test regularized variants of Method 3;
  \item rethink Method 2 around a genuinely physical matching or response-model
  framework rather than nearest-neighbour $\Delta R$ pairing.
\end{itemize}

\section{Summary}

Three cell-level photon shower-shape correction methods have been implemented
and validated on Run 3 data24 and MC23e $Z \to \ell\ell\gamma$ samples.

Method 1 is simple and robust. It preserves total energy by construction and
already brings \Reta and \Rphi to very good agreement while significantly
improving \weta. Method 3 is the best cell-level method overall. It is the only
method that explicitly matches both the mean and variance of the cell-fraction
distributions, and it clearly outperforms Method 1 for width-sensitive cases,
except in the few bins where sigma-ratio capping becomes dominant.

Method 2 required two important software fixes during this project: restoring
Francisco's original $[-1,1]$ TProfile binning and replacing the dimensionally
incorrect additive application by a proper energy shift proportional to the
cluster energy. Those fixes were necessary, but they were not sufficient. The
remaining failure of Method 2 is driven by the matching concept itself, which
introduces a systematic energy excess and forces a large renormalization that
damages the final shower-shape distributions.

The ATLAS fudge factors remain the overall benchmark. Among the cell-level
methods, Method 3 is the most promising long-term direction, while Method 1 is
the most reliable baseline. Any future revival of Method 2 will require a
fundamental redesign of the matching strategy rather than another local code
patch.

\appendix

\section{Appendix: M1 Checksums by Channel}

\begin{table}[H]
\centering
\caption{Method 1 checksums for the Zeeg channel.}
\begin{tabular}{crrc}
\toprule
Eta bin & Data events & MC events & $\sum_k \Delta_k$ \\
\midrule
0 & 7173 & 103206 & $-1.29 \times 10^{-15}$ \\
1 & 8403 & 125818 & $-3.22 \times 10^{-15}$ \\
2 & 7744 & 118984 & $5.11 \times 10^{-15}$ \\
3 & 8442 & 121045 & $2.24 \times 10^{-16}$ \\
4 & 6144 & 91277 & $4.51 \times 10^{-15}$ \\
5 & 5104 & 73833 & $7.64 \times 10^{-15}$ \\
6 & 2361 & 33585 & $-2.39 \times 10^{-15}$ \\
7 & 532 & 7993 & $-9.29 \times 10^{-16}$ \\
8 & 0 & 0 & 0 \\
9 & 1227 & 14686 & $-6.93 \times 10^{-16}$ \\
10 & 2759 & 36410 & $1.53 \times 10^{-15}$ \\
11 & 2581 & 33984 & $-4.54 \times 10^{-15}$ \\
12 & 2448 & 31577 & $5.64 \times 10^{-16}$ \\
13 & 1426 & 18497 & $-1.51 \times 10^{-15}$ \\
\bottomrule
\end{tabular}
\end{table}

\begin{table}[H]
\centering
\caption{Method 1 checksums for the Zmumug channel.}
\begin{tabular}{crrc}
\toprule
Eta bin & Data events & MC events & $\sum_k \Delta_k$ \\
\midrule
0 & 10634 & 179915 & $-7.04 \times 10^{-15}$ \\
1 & 12734 & 218187 & $4.25 \times 10^{-15}$ \\
2 & 12412 & 210938 & $9.92 \times 10^{-16}$ \\
3 & 13378 & 223845 & $4.69 \times 10^{-15}$ \\
4 & 10235 & 180080 & $8.95 \times 10^{-16}$ \\
5 & 8812 & 151930 & $-5.26 \times 10^{-15}$ \\
6 & 4119 & 69343 & $4.06 \times 10^{-15}$ \\
7 & 998 & 16827 & $-7.95 \times 10^{-16}$ \\
8 & 0 & 0 & 0 \\
9 & 2210 & 32063 & $3.35 \times 10^{-15}$ \\
10 & 5130 & 81503 & $-3.16 \times 10^{-15}$ \\
11 & 5092 & 79789 & $4.06 \times 10^{-15}$ \\
12 & 4776 & 75241 & $-4.16 \times 10^{-15}$ \\
13 & 2698 & 43472 & $6.21 \times 10^{-16}$ \\
\bottomrule
\end{tabular}
\end{table}

\section{Appendix: Method 2 Matched-Pair Statistics by Channel}

\begin{table}[H]
\centering
\caption{Method 2 matched-pair statistics for Zeeg.}
\begin{tabular}{crc}
\toprule
Eta bin & Matched pairs & $\langle |\alpha| \rangle$ for cell 38 \\
\midrule
0 & 7116 & 0.098629 \\
1 & 8393 & 0.054952 \\
2 & 7727 & 0.092285 \\
3 & 8399 & 0.090012 \\
4 & 6126 & 0.080373 \\
5 & 5085 & 0.070646 \\
6 & 2354 & 0.086366 \\
7 & 527 & 0.068719 \\
9 & 1217 & 0.094414 \\
10 & 2752 & 0.062240 \\
11 & 2570 & 0.067900 \\
12 & 2442 & 0.064379 \\
13 & 1423 & 0.068728 \\
\bottomrule
\end{tabular}
\end{table}

\begin{table}[H]
\centering
\caption{Method 2 matched-pair statistics for Zmumug.}
\begin{tabular}{crc}
\toprule
Eta bin & Matched pairs & $\langle |\alpha| \rangle$ for cell 38 \\
\midrule
0 & 10613 & 0.065505 \\
1 & 12703 & 0.073501 \\
2 & 12351 & 0.098548 \\
3 & 13334 & 0.069348 \\
4 & 10204 & 0.072185 \\
5 & 8800 & 0.063969 \\
6 & 4106 & 0.074362 \\
7 & 982 & 0.046647 \\
9 & 2205 & 0.056087 \\
10 & 5115 & 0.049852 \\
11 & 5061 & 0.048741 \\
12 & 4748 & 0.055863 \\
13 & 2701 & 0.063671 \\
\bottomrule
\end{tabular}
\end{table}

\begin{thebibliography}{9}

\bibitem{francisco}
F. Almeida, ``Shower shape corrections using cell-level reweighting'',
ATL-COM-PHYS-2021-640 (2021).

\bibitem{run3note}
M. Fernandes da Silva, ``Photon shower shape studies with Run 3 data'',
ATL-COM-PHYS-2025-662 (2025).

\bibitem{atlas-fudge}
ATLAS Collaboration, ``Photon identification and shower shape corrections'',
internal ATLAS documentation and standard reconstruction calibrations.

\end{thebibliography}

\end{document}